\begin{figure}[ht]
\centering
\begin{tikzpicture}[x=1.0cm, y=1.0cm, font=\small]

% Left panel: arrow matrix, original ordering (heavy fill-in).
% Cell (row r, col c) plots at x = (c-0.5)*0.5, y = (6-r+0.5)*0.5.
\begin{scope}[shift={(0,0)}]
\draw[thick] (0,0) rectangle (3, 3);
% diagonal
\foreach \r in {1,...,6} {
  \pgfmathsetmacro{\cx}{(\r-0.5)*0.5}
  \pgfmathsetmacro{\cy}{(6-\r+0.5)*0.5}
  \fill[engblue] ({\cx},{\cy}) circle (2.4pt);
}
% first row (columns 2..6)
\foreach \c in {2,...,6} {
  \pgfmathsetmacro{\cx}{(\c-0.5)*0.5}
  \pgfmathsetmacro{\cy}{(6-1+0.5)*0.5}
  \fill[engblue] ({\cx},{\cy}) circle (2.4pt);
}
% first column (rows 2..6)
\foreach \r in {2,...,6} {
  \pgfmathsetmacro{\cx}{(1-0.5)*0.5}
  \pgfmathsetmacro{\cy}{(6-\r+0.5)*0.5}
  \fill[engblue] ({\cx},{\cy}) circle (2.4pt);
}
% fill-in in the trailing 5x5 block (rows 2..6, cols 2..6)
\foreach \r in {2,...,6} {
  \foreach \c in {2,...,6} {
    \pgfmathsetmacro{\cx}{(\c-0.5)*0.5}
    \pgfmathsetmacro{\cy}{(6-\r+0.5)*0.5}
    \node[engred, font=\tiny] at ({\cx},{\cy}) {$\times$};
  }
}
\node[anchor=north, font=\scriptsize] at (1.5, -0.2) {pivot on hub first};
\node[anchor=south, font=\scriptsize] at (1.5, 3.15) {dense fill-in ($\times$)};
\end{scope}

% Right panel: reordered, minimal fill-in.
\begin{scope}[shift={(6.2,0)}]
\draw[thick] (0,0) rectangle (3, 3);
% diagonal
\foreach \r in {1,...,6} {
  \pgfmathsetmacro{\cx}{(\r-0.5)*0.5}
  \pgfmathsetmacro{\cy}{(6-\r+0.5)*0.5}
  \fill[engblue] ({\cx},{\cy}) circle (2.4pt);
}
% last row (columns 1..5)
\foreach \c in {1,...,5} {
  \pgfmathsetmacro{\cx}{(\c-0.5)*0.5}
  \pgfmathsetmacro{\cy}{(6-6+0.5)*0.5}
  \fill[engblue] ({\cx},{\cy}) circle (2.4pt);
}
% last column (rows 1..5)
\foreach \r in {1,...,5} {
  \pgfmathsetmacro{\cx}{(6-0.5)*0.5}
  \pgfmathsetmacro{\cy}{(6-\r+0.5)*0.5}
  \fill[engblue] ({\cx},{\cy}) circle (2.4pt);
}
\node[anchor=north, font=\scriptsize] at (1.5, -0.2) {pivot on hub last};
\node[anchor=south, font=\scriptsize] at (1.5, 3.15) {no fill-in};
\end{scope}

\node[anchor=north, font=\scriptsize, text width=11.5cm, align=center]
  at (4.6, -0.9) {
  The same arrow matrix under two orderings. Eliminating the densely
  connected hub node first (left) couples every remaining unknown and
  fills the trailing block; eliminating it last (right) leaves the
  pattern untouched. Ordering is the free lever that separates a
  tractable sparse factorisation from a dense one.
};

\end{tikzpicture}
\caption[Fill-in depends on elimination order]{Nonzero patterns of an
arrow-shaped sparse matrix under two elimination orderings. Filled
circles are original nonzeros; red crosses are fill-in, the new nonzeros
the elimination creates. Eliminating the hub node first (left panel)
couples every remaining unknown to every other and fills the trailing
block completely, so the factorisation costs and stores as if the matrix
were dense. Eliminating the hub last (right panel) creates no fill-in at
all. The reordering is a symmetric permutation
$\mat{P}\mat{A}\mat{P}^{\transpose}$ that leaves the solution unchanged
while changing the factorisation cost by orders of magnitude; a sparse
direct solver's first job is to find a good ordering, and
nested-dissection and minimum-degree heuristics are the standard tools.}
\label{fig:vol02:ch09:fillin-reordering}
\end{figure}
